\chapter{Parameter Reference}\label{chap:appB_parameters}

\providecommand{\val}[1]{\ensuremath{#1}}

This appendix is a single-place reference for the physical and numerical
parameters that define a \model{} run. Every value listed here is the
\emph{actual} default in the source code as of writing; where a parameter is a
keyword argument of a step routine, the default shown is the one that fires when
the caller passes nothing. The intent is that a reader can reproduce, or
sensibly perturb, a control integration without re-deriving constants from the
governing equations. The physics behind each group is developed in the body
chapters, which we cross-reference throughout: the ocean equations in
\ref{chap:ch04_ocean_equations}, mixing and the Gent--McWilliams parameterization
in \ref{chap:ch05_ocean_mixing}, ocean numerics in \ref{chap:ch06_ocean_numerics},
the overturning closure in \ref{chap:ch08_ocean_overturning}, the multi-level
atmosphere in \ref{chap:ch09_atmosphere_multilevel}, and the coupler and forcing
in \ref{chap:ch13_coupler_fluxes} and \ref{chap:ch14_forcing_response}.

A standing caveat: \model{} is research software, and several of these numbers
are not first-principles constants but \emph{tuning knobs} chosen for numerical
stability or for plausible large-scale behaviour at T31. The thermohaline
closure (\ref{chap:ch08_ocean_overturning}) in particular is an interim
box-model-style parameterization, and its coefficients have no observational
counterpart; they are calibrated so that the AMOC sits in a realistic range and
so that the salt-advection feedback can be pushed across a bifurcation in tipping
experiments (\ref{chap:ch17_amoc_tipping}). Read those entries as control
parameters, not measured quantities.

Symbols in the ``Code symbol'' column are the literal identifiers in the source;
the home file is given so a value can be traced and changed at its definition.

% ===================================================================
\section{Grid and resolution}
% ===================================================================

The horizontal grid is selected at import time from the environment variable
\code{CHRONOS\_RESOLUTION} (default \code{T31}); both the ocean and the
single-level atmosphere share the same regular latitude--longitude mesh, while
the active \dino{} atmosphere runs on its own Gaussian grid at the same spectral
truncation. The ocean uses a stretched vertical grid with a thin surface layer so
that the surface heat capacity is realistic and SST responds to fluxes on the
right timescale (\ref{chap:ch06_ocean_numerics}).

\begin{table}[h]
\centering
\begin{tabular}{@{}llll@{}}
\toprule
Parameter (code symbol) & Value & Units & Meaning \\
\midrule
Truncation \code{RESOLUTION} & T31 & --- & spectral truncation (default) \\
Latitudes \code{NLAT} & 48 & --- & ocean/atmos grid rows ($\approx 3.75^\circ$) \\
Longitudes \code{NLON} & 96 & --- & ocean/atmos grid columns \\
Ocean levels \code{NZ\_OCEAN} & 15 & --- & vertical $\dep$-levels \\
Total depth \code{OCEAN\_TOTAL\_DEPTH} & 5000 & m & ocean column depth \\
Stretch \code{stretch} & 1.7 & --- & vertical grid stretching exponent \\
\dino{} layers \code{layers} & 24 & --- & atmosphere $\sigma$-levels \\
\dino{} truncation \code{truncation} & T31 & --- & atmosphere spectral truncation \\
Earth radius \code{EARTH\_RADIUS}, $\Erad$ & $6.371\times10^{6}$ & m & planetary radius \\
\bottomrule
\end{tabular}
\caption{Grid and resolution. Interfaces are placed at
$\dep\,(k/\mathrm{NZ})^{1.7}$, giving a $\sim 50$~m surface layer that thickens
with depth; \code{OCEAN\_DZ} and \code{OCEAN\_DEPTH\_CENTERS} are derived from
this and stay mutually consistent.}
\label{tab:appB_grid}
\end{table}

\codref{chronos\_esm/config.py}{resolution, vertical grid, grid configs}

% ===================================================================
\section{Time steps and physical constants}
% ===================================================================

\model{} sub-cycles the ocean and atmosphere inside an hourly coupling interval
(\ref{chap:ch13_coupler_fluxes}). The single-level atmosphere uses a very short
$30$~s step (forward Euler on gravity waves, see \ref{chap:ch10_atmosphere_singlelevel});
the active \dino{} atmosphere uses an ImEx-RK3 scheme at a $20$~min step
(\ref{chap:ch09_atmosphere_multilevel}). Both can be overridden from the
environment for tuning.

\begin{table}[h]
\centering
\begin{tabular}{@{}llll@{}}
\toprule
Parameter (code symbol) & Value & Units & Meaning \\
\midrule
Ocean step \code{DT\_OCEAN} & 900 & s & ocean time step (15 min) \\
Atmos step \code{DT\_ATMOS} & 30 & s & single-level atmosphere step \\
\dino{} step \code{dt\_minutes} & 20 & min & multi-level atmosphere step \\
Coupling \code{COUPLING\_INTERVAL} & 3600 & s & ocean--atmos exchange interval \\
CG tol.\ \code{CG\_TOL} & $10^{-6}$ & --- & conjugate-gradient tolerance \\
CG max iter \code{CG\_MAX\_ITER} & 1000 & --- & CG iteration cap \\
Softplus sharpness \code{EPSILON\_SMOOTH} & $10^{-3}$ & --- & $\softplus$ scale (precip) \\
Gravity \code{GRAVITY}, $\grav$ & 9.81 & m\,s$^{-2}$ & gravitational acceleration \\
Rotation \code{OMEGA}, $\Om$ & $7.2921\times10^{-5}$ & rad\,s$^{-1}$ & Earth angular velocity \\
Seawater density \code{RHO\_WATER}, $\densz$ & 1025 & kg\,m$^{-3}$ & reference seawater density \\
Air density \code{RHO\_AIR} & 1.225 & kg\,m$^{-3}$ & reference air density \\
$c_{p,\mathrm{air}}$ \code{CP\_AIR} & 1004 & J\,kg$^{-1}$\,K$^{-1}$ & specific heat of dry air \\
$c_{p,\mathrm{sea}}$ \code{CP\_WATER} & 3985 & J\,kg$^{-1}$\,K$^{-1}$ & specific heat of seawater \\
$R_{\mathrm{air}}$ \code{R\_AIR} & 287 & J\,kg$^{-1}$\,K$^{-1}$ & gas constant, dry air \\
$L_v$ \code{LATENT\_HEAT\_VAPORIZATION} & $2.5\times10^{6}$ & J\,kg$^{-1}$ & latent heat of vaporization \\
$\sigma$ \code{STEFAN\_BOLTZMANN} & $5.67\times10^{-8}$ & W\,m$^{-2}$\,K$^{-4}$ & Stefan--Boltzmann constant \\
$p_0$ \code{P0} & $10^{5}$ & Pa & reference pressure \\
\bottomrule
\end{tabular}
\caption{Time steps, solver tolerances, and physical constants.
$\kappa = R_{\mathrm{air}}/c_{p,\mathrm{air}}$ is derived. Note that the
hourly coupling interval is sub-cycled by both the ocean ($900$~s) and the
atmosphere.}
\label{tab:appB_time}
\end{table}

\codref{chronos\_esm/config.py}{time stepping, constants, tuning coefficients}

% ===================================================================
\section{Ocean physics}
% ===================================================================

These are the defaults of \code{step\_ocean} (the ocean stepper). They set the
viscous and diffusive closure (\ref{chap:ch05_ocean_mixing}), the Gent--McWilliams
eddy-bolus coefficient \citep{gentmcwilliams1990}, the bottom drag entering the
geostrophic balance, and the safety clamps that keep the integration physical
(\ref{chap:ch06_ocean_numerics}). The equation of state is the linear form
$\dens = \densz - \alpha(\Temp - \Temp_0) + \beta(\Salt - \Salt_0)$ used
throughout \ref{chap:ch04_ocean_equations}.

\begin{table}[h]
\centering
\begin{tabular}{@{}llll@{}}
\toprule
Parameter (code symbol) & Value & Units & Meaning \\
\midrule
Lateral viscosity \code{Ah} & $2.0\times10^{5}$ & m$^{2}$\,s$^{-1}$ & harmonic momentum viscosity \\
Biharmonic visc.\ \code{Ab} & 0 & m$^{4}$\,s$^{-1}$ & biharmonic viscosity (off) \\
GM coefficient \code{kappa\_gm} & 1000 & m$^{2}$\,s$^{-1}$ & Gent--McWilliams bolus \\
Tracer diff.\ \code{kappa\_h} & 500 & m$^{2}$\,s$^{-1}$ & interior horizontal diffusivity \\
Biharmonic tracer \code{kappa\_bi} & 0 & m$^{4}$\,s$^{-1}$ & biharmonic tracer diff.\ (off) \\
Smagorinsky \code{smag\_constant} & 0.1 & --- & dynamic-viscosity constant \\
Bottom drag \code{r\_drag} & $5.0\times10^{-2}$ & s$^{-1}$ & linear drag in geostrophy \\
Barotropic drag \code{r\_drag\_bt} & $1/(86400\cdot25)$ & s$^{-1}$ & Stommel 25-day drag \\
Shapiro filter \code{shapiro\_strength} & 0 & --- & $2\Delta x$ noise cleanup (off) \\
Polar diffusivity (interior) & 20000 & m$^{2}$\,s$^{-1}$ & enhanced diff.\ at $|\,\lat\,|$ poles \\
EOS $\alpha$ & 0.2 & kg\,m$^{-3}$\,K$^{-1}$ & thermal expansion \\
EOS $\beta$ & 0.8 & kg\,m$^{-3}$\,psu$^{-1}$ & haline contraction \\
EOS $\Temp_0$ & 283.15 & K & reference temperature \\
EOS $\Salt_0$ & 35.0 & psu & reference salinity \\
Salinity clip & $[30,\,38]$ & psu & hard stability clamp \\
Temperature clip & $[250,\,320]$ & K & hard stability clamp \\
Freezing floor & 271.35 & K & $\approx-1.8^\circ$C sea-ice stub \\
Sponge salinity \code{S\_REF} & 35.0 & psu & high-lat sponge target \\
Sponge timescale \code{tau\_sponge} & 180 & days & polar ($>65^\circ$) relaxation \\
Global relax \code{tau\_global} & 3 & yr & weak global salinity relaxation \\
Global mean \code{S\_REF\_GLOBAL} & 34.7 & psu & conserved volume-mean salinity \\
\bottomrule
\end{tabular}
\caption{Ocean physics defaults from \code{step\_ocean}. The two-tier salinity
treatment (a strong $6$-month polar sponge plus a weak $3$-year global
relaxation, with the volume mean renormalized each step) is described in
\ref{chap:ch07_ocean_prognostic_core}; the clamps are stability nets and are not
hit in a healthy control run.}
\label{tab:appB_ocean}
\end{table}

\codref{chronos\_esm/ocean/veros\_driver.py}{step\_ocean defaults, EOS, sponges, clamps}

% ===================================================================
\section{Thermohaline-closure knobs}
% ===================================================================

The thermohaline overturning closure
(\code{thc\_overturning\_velocity}, \ref{chap:ch08_ocean_overturning}) adds a
depth-integral-zero Atlantic meridional cell whose strength scales as
$\code{thc\_k\_vel}\cdot\softplus(\text{contrast}/\code{drho\_scale})$, where the
contrast is the subpolar-minus-subtropical upper-ocean density difference,
optionally haline-amplified. Because the cell carries no net transport it
survives the barotropic corrector and gives density a clean pathway to the
$26.5^\circ$N overturning, so that $\partial(\mathrm{AMOC})/\partial(\dens)$ is
nonzero and sign-correct. Raising \code{thc\_haline\_gain} above unity amplifies
the salt-advection feedback past a loop gain of one, making the AMOC bistable for
the tipping experiments in \ref{chap:ch17_amoc_tipping}
\citep{stommel1961, rahmstorf2005, marotzke1991}. These are control knobs, not
observed constants.

\begin{table}[h]
\centering
\begin{tabular}{@{}llll@{}}
\toprule
Parameter (code symbol) & Value & Units & Meaning \\
\midrule
Velocity scale \code{thc\_k\_vel} & $1.0\times10^{-4}$ & m\,s$^{-1}$ & overturning strength scale \\
Haline gain \code{thc\_haline\_gain} & 1.0 & --- & salt-feedback amplification ($1=$ off) \\
Contrast depth \code{thc\_contrast\_depth\_m} & None ($\to$ \code{upper\_depth\_m}) & m & layer for density contrast \\
Limb depth \code{upper\_depth\_m} & 1100 & m & overturning-cell upper-limb depth \\
Contrast scale \code{drho\_scale} & 0.1 & kg\,m$^{-3}$ & $\softplus$ contrast normalization \\
Haline coeff.\ \code{BETA\_S} & 0.78 & kg\,m$^{-3}$\,psu$^{-1}$ & $\partial\dens/\partial\Salt$ for amplification \\
Subpolar band \code{subpolar} & $(45,\,65)$ & $^\circ$N & dense-water-formation band \\
Subtropical band \code{subtropical} & $(10,\,30)$ & $^\circ$N & reference band \\
Hosing \code{hosing\_sv} & 0 & \unskip$\Sv$ & subpolar freshwater forcing \\
Hosing band \code{band} & $(45,\,65)$ & $^\circ$N & hosing footprint \\
\bottomrule
\end{tabular}
\caption{Thermohaline-closure and AMOC-hosing knobs. \code{thc\_k\_vel=0}
disables the closure entirely. A shallow \code{thc\_contrast\_depth\_m}
($\sim300$~m) gives a surface freshwater hosing more leverage on deep-water
formation than the default $0$--$1100$~m mean; see
\ref{chap:ch17_amoc_tipping}.}
\label{tab:appB_thc}
\end{table}

\codref{chronos\_esm/ocean/overturning.py}{thc\_overturning\_velocity, subpolar\_hosing\_salt\_tendency}

% ===================================================================
\section{Atmosphere}
% ===================================================================

The active atmosphere is the \dino{} differentiable spectral primitive-equation
dycore (\ref{chap:ch09_atmosphere_multilevel}), run with Held--Suarez-style
forcing \citep{heldsuarez1994} and the architecture behind NeuralGCM
\citep{kochkov2024neuralgcm}. The defaults below are the \code{DinoAtmosphere}
constructor arguments. The eddy-permitting tuning is deliberate: the
hyperdiffusion timescale is weak enough ($6$~h) that T31 baroclinic eddies can
grow, and a small rotational wind seed breaks the exactly axisymmetric
equilibrium so that mid-latitude surface westerlies are earned dynamically
(\ref{chap:ch18_jets_stormtracks}). The legacy single-level path
(\ref{chap:ch10_atmosphere_singlelevel}) is retained but is not the active
atmosphere.

\begin{table}[h]
\centering
\begin{tabular}{@{}llll@{}}
\toprule
Parameter (code symbol) & Value & Units & Meaning \\
\midrule
Diffusion timescale \code{diffusion\_tau\_hours} & 6.0 & h & hyperdiffusion damping time \\
Diffusion order \code{diffusion\_order} & 2 & --- & $\lap$-order of the filter \\
Boundary-layer top \code{sigma\_b} & 0.7 & --- & $\sigma$ of friction layer top \\
Friction rate \code{kf\_per\_day} & 1.0 & day$^{-1}$ & boundary-layer Rayleigh drag \\
Free-trop.\ relax \code{ka\_per\_day} & $1/40$ & day$^{-1}$ & Newtonian cooling rate \\
Surface relax \code{ks\_per\_day} & $1/4$ & day$^{-1}$ & near-surface relaxation rate \\
Min.\ equil.\ temp.\ \code{minT} & 200 & K & stratospheric $\Temp_{eq}$ floor \\
Horizontal $\Delta\Temp$ \code{dThz} & 10 & K & equator--pole $\Temp_{eq}$ contrast \\
Condensation time \code{tau\_cond\_hours} & 3.0 & h & large-scale precip relaxation \\
Initial RH \code{init\_rh} & 0.5 & --- & initial relative humidity \\
Seed wind \code{seed\_wind\_ms} & 2.0 & m\,s$^{-1}$ & symmetry-breaking vorticity seed \\
Orography \code{orography} & True & --- & ETOPO topography on/off \\
\bottomrule
\end{tabular}
\caption{\dino{} atmosphere defaults (\code{DinoAtmosphere.\_\_init\_\_}).
The time step and truncation appear in Tables~\ref{tab:appB_grid} and
\ref{tab:appB_time}.}
\label{tab:appB_atmos}
\end{table}

\codref{chronos\_esm/atmos/dino\_atmos.py}{DinoAtmosphere constructor, forcing, seed}

% ===================================================================
\section{Coupler and forcing}
% ===================================================================

The coupler (\ref{chap:ch13_coupler_fluxes}) computes bulk surface fluxes
\citep{largeyeager2004}, applies an SST flux correction, and injects external
radiative forcing. The CO$_2$
forcing follows the logarithmic Myhre form
$F = 5.35\,\ln(C/C_0)$~W\,m$^{-2}$ \citep{myhre1998}, zero at the preindustrial
reference $C_0 = 280$~ppm, and is the channel for the forcing-response and CMIP-style
experiments (\ref{chap:ch14_forcing_response}, \ref{chap:ch20_cmip7})
\citep{eyring2016cmip6}. The flux correction has two modes: a strong Haney
restoring of SST to the WOA target ($\tau\approx30$~days, control mode), or a
frozen q-flux with weak anomaly restoring ($\tau$ long, free / forcing-responsive
mode) \citep{haney1971}.

\begin{table}[h]
\centering
\begin{tabular}{@{}llll@{}}
\toprule
Parameter (code symbol) & Value & Units & Meaning \\
\midrule
CO$_2$ coefficient & 5.35 & W\,m$^{-2}$ & Myhre log-forcing coefficient \\
CO$_2$ reference \code{co2\_ref} & 280 & ppm & preindustrial reference (zero forcing) \\
Restoring time \code{RESTORE\_TAU\_DAYS} & 30 & days & control-mode SST restoring \\
Free-mode restoring \code{restore\_tau\_days} & 3650 & days & q-flux anomaly restoring \\
Drag coeff.\ \code{CD} & $1.3\times10^{-3}$ & --- & bulk wind-stress drag \\
Coupler air density \code{RHO\_AIR} & 1.2 & kg\,m$^{-3}$ & bulk-flux air density \\
Ocean drag \code{DRAG\_COEFF\_OCEAN} & $1.2\times10^{-3}$ & --- & surface drag over ocean \\
Land drag \code{DRAG\_COEFF\_LAND} & $2.0\times10^{-3}$ & --- & surface drag over land \\
Ocean albedo \code{ALBEDO\_OCEAN} & 0.07 & --- & shortwave albedo, ocean \\
Land albedo \code{ALBEDO\_LAND} & 0.30 & --- & shortwave albedo, land \\
Land emissivity \code{EMISSIVITY\_LAND} & 0.96 & --- & longwave emissivity, land \\
Wind-stress clip & $\pm0.3$ & Pa & $\tau_x,\tau_y$ safety clamp \\
Net-heat clip & $\pm1500$ & W\,m$^{-2}$ & surface heat-flux clamp \\
Coupling interval \code{interval} & 1.0 & day & coupled-step interval \\
\bottomrule
\end{tabular}
\caption{Coupler and forcing parameters. The two restoring timescales select the
flux-correction mode; in control runs the model restores SST strongly to WOA
($30$~days), while forcing-response runs freeze a diagnosed q-flux and restore
only anomalies on a multi-year timescale.}
\label{tab:appB_coupler}
\end{table}

\codref{chronos\_esm/coupler/dino\_coupling.py}{co2\_forcing\_wm2, ocean\_fluxes\_jax, restoring}
\codref{chronos\_esm/coupler/dino\_step.py}{DinoCoupledModel: interval, thc/flux-correction wiring}

% ===================================================================
\section{A note on provenance and reproducibility}
% ===================================================================

Two practical points close this reference. First, the resolution, both
atmospheric time steps, and the contrast depth can be overridden at run time
(\code{CHRONOS\_RESOLUTION}, \code{CHRONOS\_DT\_ATMOS}, and constructor keywords),
so a reproduced run should record the environment and the explicit keyword
arguments, not merely cite this table. Second, several values --- the THC knobs,
the salinity sponge and global-relaxation timescales, the ocean clamps, and the
hyperdiffusion time --- are pragmatic choices that trade physical fidelity against
the stability of a low-resolution, partly diagnostic ocean. They are documented
here so that they can be changed honestly and with their consequences understood,
which is the point of running a fully differentiable model
(\ref{chap:ch03_differentiable}, \ref{chap:ch15_differentiable_applications})
\citep{jax2018, vallis2017, griffies2004, semtner1976}.
